% 熱物理学 熱物理学-3 問題7（022-genshi） 解いた日 2026-11-12
%% kotae
\toi{1}{◎}{}
\[ \frac{\partial}{\partial V}V^{2}=2V, \qquad \frac{\partial}{\partial T}(2TV)=2V \]
より条件2を満たす。
\toi{2}{◎}{}
\begin{align*}
F(T,V)&=\int_{T^{*}}^{T}V^{*2}\,\dd T'+\int_{V^{*}}^{V}2TV'\,\dd V'\\
&=\bigl[V^{*2}T'\bigr]_{T^{*}}^{T}+\bigl[TV'^{2}\bigr]_{V^{*}}^{V}\\
&=(V^{*2}T-V^{*2}T^{*})+(TV^{2}-TV^{*2})\\
&=TV^{2}-T^{*}V^{*2}
\end{align*}
\[ \left(\pdv{F}{T}\right)_{V}=V^{2}, \qquad \left(\pdv{F}{V}\right)_{T}=2TV \]
\toi{3}{○}{記述}
$T=T_{0}(1+s)$，$V=V_{0}(1-s/2)$，$\dd T=T_{0}\,\dd s$，$\dd V=-\frac{1}{2}V_{0}\,\dd s$
\begin{align*}
\ayamari{F(T,V)}{I_{C}}&=\int_{C}\omega=\int_{C}V^{2}\,\dd T+2TV\,\dd V\\
&=\int_{0}^{1}V_{0}^{2}\left(1-\frac{s}{2}\right)^{2}T_{0}\,\dd s\\
&\qquad+\int_{0}^{1}2T_{0}V_{0}(1+s)\left(1-\frac{s}{2}\right)\\
&\qquad\qquad\cdot\left(-\frac{1}{2}V_{0}\right)\dd s\\
&=T_{0}V_{0}^{2}\int_{0}^{1}\biggl(\left(1-\frac{s}{2}\right)^{2}\\
&\qquad-(1+s)\left(1-\frac{s}{2}\right)\biggr)\dd s\\
&=T_{0}V_{0}^{2}\left[-\frac{3}{4}s^{2}+\frac{1}{4}s^{3}\right]_{0}^{1}=-\frac{1}{2}T_{0}V_{0}^{2}
\end{align*}
\naze{線積分 $\int_{C}\omega$ は曲線 $C$ で決まる1つの数で，$(T,V)$ の関数 $F$ ではない（等しいのは $F(\mathrm{Q})-F(\mathrm{P})$）。$I_{C}$ のように別の名前を付ける。}
\begin{align*}
F(\mathrm{Q})-F(\mathrm{P})&=2T_{0}\left(\frac{V_{0}}{2}\right)^{2}-T_{0}V_{0}^{2}\\
&=-\frac{1}{2}T_{0}V_{0}^{2}
\end{align*}
一致。
\naze{両者が等しくなる理由が無い。$\omega$ は (1) で条件2を確かめた全微分なので，線積分は原始関数の終点と始点の値の差 $F(\mathrm{Q})-F(\mathrm{P})$ になる（基準点による定数 $T^{*}V^{*2}$ は差で消える）。}
%% kizuki
\item 線積分 $\int_{C}\omega$ が $F(\text{終点})-F(\text{始点})$ に等しいのは，$\omega$ が全微分（原始関数をもつ）ときだけである。(1) で確かめた条件2を理由として書き，線積分には $I_{C}$ など原始関数 $F$ と別の名前を付ける。
\item (2) では最初に第2項の $T$ を $T^{*}$ にしかけて自分で直した。原始関数の構成の第2項 $\int_{V^{*}}^{V}B(T,V')\,\dd V'$ の $T$ は終点の $T$（固定した値）である。
